\section{Discussion and Conclusion}

The upgraded WebCoEvo evidence changes the paper's center of gravity. The earlier version treated the final benchmark as a narrower website-drift ablation with unresolved held-out access evidence. The current result is cleaner: the same standalone runner reports complete 68-row training and 167-row held-out validation comparisons across \linkding{} V1.45.0, website V2, and website V3. That makes the empirical claim both stronger and more constrained.

The strongest methodological point is trace-audited separation of rule layers. \expel{} rules and reflection rules are not treated as one opaque prompt. The runner records which task-experience rules were injected, which reflection rules were injected, which rulebook path was selected, and whether failures occurred during reset or during the agent rollout. This evidence standard matters because a reflection-rule paper can otherwise confuse a rulebook existing on disk with a rule actually reaching the model.

The strongest empirical point is that website drift changes which kind of rule is useful. On the training tasks from \linkding{} V1.45.0, \expel{} rules are extremely helpful. On the training tasks under website V3, the same class of task-experience rules collapses to \(8/68=11.8\%\), while R1 reflection rules reach \(65/68=95.6\%\). On the held-out validation tasks under website V3, R1 and R4 tie at \(97/167=58.1\%\), improving over \expel{}-style rules at \(66/167=39.5\%\). The lesson is not that task-experience rules are bad. It is that task-experience rules and reflection rules solve different parts of the adaptation problem.

\paragraph{Rule refinement.}
The natural next rulebook should start from R1, not simply from the latest merge. R1 remains the strongest rule set on the training tasks under website V3 and ties the best held-out validation score. R4 contributes useful evidence about login-next preservation and finalization, but its training-task regression means those edits need a narrower promotion gate. A good next version should preserve R1's process and functional behavior while selectively borrowing R4's access, surface, content, and runtime gains.

\paragraph{Environment hardening.}
The environment generator should remain capability-gated. Saturated drift types should be hardened by breaking one primary assumption while preserving task semantics and evaluator validity. Frontier drift types should be held steady for rule validation. Infrastructure failures should trigger runner or evaluator repair, not rule mining. This keeps the co-evolution loop honest: the next website should expose an agent limitation, not manufacture an invalid task.

\paragraph{Cross-site transfer.}
The open claim is whether WebCoEvo is a \linkding{} repair procedure or a reusable web-agent adaptation protocol. The next transfer step should use a second local web application with deterministic reset, credentials, seed data, and stable evaluators. We would map the same seven drift categories into that site, run \expel{}-style versus compact reflection-rule baselines, and then test whether the mined rules remain general or become site-specific.

\paragraph{Limitations.}
The current evidence remains single-site. The website versions are generated from one application family, and the strongest held-out validation score under website V3 is still only \(58.1\%\). Some rule edits improve one drift family while regressing another, and aggregate scores can hide those tradeoffs. Finally, the method depends on trace quality: if the runner fails to expose reset errors, rule-selection misses, or prompt-injection state, the co-evolution loop can mine the wrong evidence.

In conclusion, \webcoevo{} turns web-agent website drift from a static benchmark into an auditable improvement loop. It profiles failures, generates capability-targeted \linkding{} website versions, separates task-experience and reflection rules, and shows that compact reflection rules recover substantial performance under website V3 drift. The current best evidence is R1: \(65/68=95.6\%\) on the training tasks under website V3 and \(97/167=58.1\%\) on the held-out validation tasks under website V3. The remaining work is concrete: preserve R1's robust behaviors, repair the drift-family regressions exposed by R4, and test the same protocol beyond \linkding{}.
